The @@M0@@-ring @@M1@@ is characterized by the following universal property:

\sgaStableTarget{sga6.e05.prop.4-2}\textbf{Proposition 4.2.} \uline{For every} @@M2@@-@@M3@@-\uline{algebra} @@M4@@, \uline{and every family} @@M5@@ \uline{of elements of} @@M6@@, \uline{there exists exactly one} @@M7@@-@@M8@@-\uline{homomorphism} @@M9@@ \uline{from} @@M10@@ \uline{to} @@M11@@ \uline{such that for every} @@M12@@, @@M13@@.

For @@M14@@ to be a @@M15@@-homomorphism, we must have @@M16@@ for all @@M17@@; this defines @@M18@@ uniquely. The homomorphism thus obtained is indeed a @@M19@@-homomorphism, since to verify that @@M20@@ it suffices to verify it on a system of generators of @@M21@@ over @@M22@@, which follows from the definition of @@M23@@ and relations (2.4.2).

\sgaStableTarget{sga6.e05.para.4-3}4.3. We can make the same construction using the @@M24@@ in place of the @@M25@@. Let @@M26@@; to define the @@M27@@-ring structure on @@M28@@, we specify a ring homomorphism: @@M29@@, by:
\sgaStableTarget{sga6.e05.eq.4-3-1}@@M30@@
We then verify the commutativity of (3.7.2), which shows that we do obtain a @@M31@@-ring structure.

The @@M32@@-ring @@M33@@ also satisfies the universal property 4.2; it is therefore canonically isomorphic (as a @@M34@@-ring) to the @@M35@@-ring @@M36@@.

\sgaStableTarget{sga6.e05.def.4-4}\textbf{Definition 4.4.} \uline{The} @@M37@@-\uline{ring} @@M38@@ \uline{will be called the} free @@M39@@-ring generated over @@M40@@ by the family of generators @@M41@@.

\sgaStableTarget{sga6.e05.def.4-5}\textbf{Definition 4.5.} \uline{Let} @@M42@@ \uline{be a} @@M43@@-\uline{ring}, @@M44@@ \uline{an ideal of} @@M45@@. \uline{We say that} @@M46@@ \uline{is} \uline{a} @@M47@@-ideal of @@M48@@ \uline{if for every} @@M49@@, \uline{and every} @@M50@@, @@M51@@.

If @@M52@@ is a @@M53@@-homomorphism, the kernel of @@M54@@ is a @@M55@@-ideal.

Let @@M56@@ be a @@M57@@-ideal; the composite homomorphism: @@M58@@ vanishes on @@M59@@ and therefore factors through @@M60@@, which is thus endowed
